\subsection{Memory System Effectiveness}
\label{sec:exp_memory}

We evaluate the memory system along four axes: whether experience accumulation yields measurable efficiency gains across sessions (\S\ref{sec:exp_mem_eff}), whether condensed memory outperforms verbose alternatives (\S\ref{sec:exp_mem_ablation}), whether the hierarchical architecture supports competitive long-term fact retention without vector databases (\S\ref{sec:exp_mem_fact}), and whether context size remains bounded as skills and memory scale (\S\ref{sec:exp_mem_explosion}).


\subsubsection{Cross-Session Efficiency Gains}
\label{sec:exp_mem_eff}

We compare GenericAgent against CodeX, Claude Code, and OpenClaw on a repeated-task protocol: each system downloads five different HuggingFace datasets in five independent sessions (fresh context per session).
We record \emph{operation time} (wall-clock time minus download I/O) and \emph{token consumption} per session.
All systems share the same backbone LLM and environment.

As shown in Table~\ref{tab:memory_improvement}, the three baselines exhibit flat trajectories---CodeX fluctuates between 92--98\,s, Claude Code between 88--93\,s, and OpenClaw between 108--115\,s, with token counts similarly stable.
GenericAgent starts at a comparable level (102\,s, 200k tokens) but improves monotonically: by the fifth session, operation time drops to 65\,s ($-$39.2\%) and token consumption to 99k ($-$50.4\%).
The improvement is driven by Reflect-mode distillation: after each session the system extracts reusable patterns (API construction, authentication flow, error recovery) into L3 SOPs, which subsequent sessions load and follow directly rather than re-deriving from scratch.

\begin{table}[t]
\centering
\caption{Repeated HuggingFace download task across five sessions. GenericAgent improves continuously; baselines remain flat.}
\label{tab:memory_improvement}
\begin{tabular}{lcccccc}
\toprule
\textbf{System} & \textbf{S1} & \textbf{S2} & \textbf{S3} & \textbf{S4} & \textbf{S5} & \textbf{$\Delta$ S1$\to$S5} \\
\midrule
\multicolumn{7}{l}{\emph{Operation Time (s)}} \\
\midrule
CodeX        & 95  & 98  & 92  & 97  & 94  & $-$1.1\% \\
Claude Code  & 88  & 91  & 89  & 93  & 90  & +2.3\% \\
OpenClaw     & 112 & 108 & 115 & 110 & 113 & +0.9\% \\
GenericAgent & 102 & 78  & 71  & 62  & 65  & \textbf{$-$39.2\%} \\
\midrule
\multicolumn{7}{l}{\emph{Token Consumption ($\times 10^3$)}} \\
\midrule
CodeX        & 185 & 189 & 182 & 190 & 187 & +0.8\% \\
Claude Code  & 173 & 175 & 171 & 178 & 174 & +0.8\% \\
OpenClaw     & 216 & 219 & 212 & 220 & 218 & +1.0\% \\
GenericAgent & 200 & 145 & 129 & 112 & 99  & \textbf{$-$50.4\%} \\
\bottomrule
\end{tabular}
\end{table}


\subsubsection{Memory Content Ablation}
\label{sec:exp_mem_ablation}

To isolate the effect of memory \emph{content selection}, we fix the backbone (GPT-5.4), framework, and task set, varying only the memory injected into the prompt.
The benchmark is SOP-Bench dangerous\_goods, a rule-heavy procedural task requiring input validation, score aggregation, and risk-level mapping---rules the model cannot reliably infer from the task description alone.
Four configurations are compared:
\textbf{No-Memory} (0 tokens),
\textbf{Condensed} (165 tokens, critical rules only),
\textbf{Full} (575 tokens, complete SOP),
and \textbf{Redundant} (288 tokens, condensed + background explanations).

Results appear in Table~\ref{tab:memory_ablation}.
Without memory the model achieves only 13.87\% Task Success Rate (TSR), confirming that external procedural knowledge is necessary.
Full-Memory improves to 52.44\%, yet Condensed-Memory reaches 66.48\%---14 points higher at less than one-third the token cost.
Redundant-Memory matches Condensed at 66.48\% but consumes 1.7$\times$ more tokens, indicating that the additional explanatory content contributes no marginal benefit.
The takeaway is that \emph{higher information density, not larger memory volume}, drives performance: retaining only the rules that directly alter model behavior yields both lower cost and higher accuracy.

\begin{table}[t]
\centering
\caption{Memory content ablation on SOP-Bench dangerous\_goods (GPT-5.4). Condensed memory achieves the highest TSR at the lowest token cost.}
\label{tab:memory_ablation}
\begin{tabular}{lcc}
\toprule
\textbf{Configuration} & \textbf{Size (tokens)} & \textbf{TSR (\%) $\uparrow$} \\
\midrule
No-Memory   & 0   & 13.87 \\
Full        & 575 & 52.44 \\
Redundant   & 288 & 66.48 \\
Condensed   & 165 & \textbf{66.48} \\
\bottomrule
\end{tabular}
\end{table}


\subsubsection{Long-Term Fact Retention}
\label{sec:exp_mem_fact}

Most long-term memory systems rely on embedding models and vector stores for retrieval.
GenericAgent instead uses hierarchical indexing (L1 metadata $\to$ L2 facts $\to$ L3 SOPs) with action-verified updates and no external retrieval infrastructure.
We benchmark on Locomo10 (Category 5 excluded), covering Multi-Hop, Temporal, Open-Domain, and Single-Hop fact recall.
Baselines include two embedding-based systems---\textbf{Mem0} and \textbf{A-MEM} (both with text-embedding-3-small)---and \textbf{OpenClaw} (no embeddings).
All use GPT-5.4.

Table~\ref{tab:long_term_memory} reports F1 and BLEU-1.
GenericAgent leads in every category despite using no embedding model.
The margin is largest on Multi-Hop (F1 43.33 vs.\ 39.32 for Mem0) and Temporal (52.23 vs.\ 50.03), where cross-fact reasoning benefits from GA's topic-based co-location of related memories.
OpenClaw, the other embedding-free system, trails substantially (F1 9.56--23.44), confirming that the advantage comes from GA's memory \emph{organization}---hierarchical distillation and action-verified quality control---rather than from simply avoiding embeddings.

\begin{table}[t]
\centering
\caption{Long-term fact retention on Locomo10 (GPT-5.4). GenericAgent outperforms embedding-based baselines without vector retrieval.}
\label{tab:long_term_memory}
\begin{tabular}{lcccccccc}
\toprule
& \multicolumn{2}{c}{\textbf{Multi-Hop}} & \multicolumn{2}{c}{\textbf{Temporal}} & \multicolumn{2}{c}{\textbf{Open-Domain}} & \multicolumn{2}{c}{\textbf{Single-Hop}} \\
\cmidrule(lr){2-3} \cmidrule(lr){4-5} \cmidrule(lr){6-7} \cmidrule(lr){8-9}
& F1 & BLEU & F1 & BLEU & F1 & BLEU & F1 & BLEU \\
\midrule
Mem0         & 39.32 & 28.43 & 50.03 & 43.33 & 18.32 & 13.80 & 40.32 & 32.33 \\
A-MEM        & 29.03 & 20.48 & 46.83 & 38.84 & 13.11 & 12.94 & 44.68 & 37.01 \\
OpenClaw     & 21.43 & 18.41 & 22.56 & 20.43 &  9.56 &  9.03 & 23.44 & 24.21 \\
GenericAgent & \textbf{43.33} & \textbf{39.96} & \textbf{52.23} & \textbf{51.11} & \textbf{20.41} & \textbf{15.31} & \textbf{45.69} & \textbf{40.66} \\
\bottomrule
\end{tabular}
\end{table}


\subsubsection{Context Explosion Prevention}
\label{sec:exp_mem_explosion}

A practical concern for long-running agents is unbounded prompt growth: as skills and memory accumulate, does the context window explode?
We stress-test this by installing 20 skills in each system, running diverse tasks over an extended period, then measuring the full prompt length on a trivial input (``Hello'')---a query that should require near-zero context beyond the system prompt and tool definitions.

Table~\ref{tab:context_explosion} shows that baseline prompts range from 22.8k (Claude Code) to 43.3k tokens (OpenClaw), whereas GenericAgent requires only 2,298 tokens---a 90--95\% reduction.
The gap arises because baselines eagerly load skill descriptions, accumulated metadata, and historical context regardless of relevance.
GenericAgent's L1 index evaluates relevance \emph{before} loading: for ``Hello'', no L2/L3 content qualifies, so none enters the prompt; the 20 skill implementations remain on disk until explicitly invoked.
Combined with the compact eight-tool schema, this keeps the base prompt small and stable, ensuring that long-term capability growth does not degrade per-turn reasoning quality.

\begin{table}[t]
\centering
\caption{Full prompt length on a trivial input after installing 20 skills and extended usage. GenericAgent prevents context explosion.}
\label{tab:context_explosion}
\begin{tabular}{lc}
\toprule
\textbf{System} & \textbf{Prompt Length (tokens)} \\
\midrule
Claude Code  & 22,821 \\
CodeX        & 23,932 \\
OpenClaw     & 43,321 \\
GenericAgent & \textbf{2,298} \\
\bottomrule
\end{tabular}
\end{table}
